\documentclass[12pt, letterpaper]{article}

\usepackage[utf8]{inputenc}
\usepackage[T1]{fontenc}
\usepackage[spanish]{babel}
\usepackage{geometry}
\usepackage{amsmath, amssymb, amsthm}
\usepackage{booktabs}
\usepackage{array}
\usepackage{microtype}
\usepackage{setspace}
\usepackage{parskip}
\usepackage{titlesec}
\usepackage{fancyhdr}
\usepackage{enumitem}
\usepackage{bm}
\usepackage{mathtools}
\usepackage{tcolorbox}
\usepackage{hyperref}

\tcbuselibrary{skins, breakable}

\geometry{top=2.8cm, bottom=3.0cm, left=3.2cm, right=3.2cm}
\setstretch{1.20}

\pagestyle{fancy}
\fancyhf{}
\fancyhead[C]{\small\textsc{Econometria · Gujarati --- Ejercicio 11.5}}
\fancyfoot[C]{\small\thepage}
\renewcommand{\headrulewidth}{0.4pt}
\renewcommand{\footrulewidth}{0pt}

\titleformat{\section}
  {\large\bfseries\scshape}{\thesection.}{0.6em}{}
  [\vspace{-4pt}\rule{\linewidth}{0.4pt}]
\titleformat{\subsection}{\normalsize\bfseries}{\thesubsection.}{0.5em}{}
\titleformat{\subsubsection}{\normalsize\itshape}{}{0em}{}

\newtheorem{prop}{Proposicion}
\newtheorem{lema}{Lema}
\theoremstyle{remark}
\newtheorem{obs}{Observacion}

\newtcolorbox{clave}{
  colback=white, colframe=black,
  boxrule=0.5pt, arc=3pt,
  left=8pt, right=8pt, top=6pt, bottom=6pt,
  breakable
}

\newcommand{\E}{\mathbb{E}}
\newcommand{\Var}{\operatorname{Var}}
\newcommand{\sw}{\sum w_i}
\newcommand{\swx}{\sum w_i X_i}
\newcommand{\swy}{\sum w_i Y_i}
\newcommand{\swxy}{\sum w_i X_i Y_i}
\newcommand{\swxx}{\sum w_i X_i^2}
\newcommand{\Ybar}{\bar{Y}^*}
\newcommand{\Xbar}{\bar{X}^*}
\newcommand{\ystar}{y_i^*}
\newcommand{\xstar}{x_i^*}

\begin{document}

\begin{center}
  {\LARGE\bfseries Minimos Cuadrados Ponderados:}\\[6pt]
  {\LARGE\bfseries Representacion en Desviaciones respecto}\\[6pt]
  {\LARGE\bfseries a las Medias Ponderadas}\\[12pt]
  {\large\itshape Demostracion completa de las ecuaciones (11.3.8) y (11.3.9)}\\[14pt]
  \rule{0.55\linewidth}{0.6pt}\\[8pt]
  {\normalsize Ejercicio~11.5 --- \textit{Econometria}, Gujarati \& Porter, 5.a ed.}\\[4pt]
  {\small\textsc{Con manipulacion algebraica paso a paso, lemas auxiliares
  y verificacion de la varianza}}
\end{center}

\vspace{10pt}

\noindent\textbf{Planteamiento.}\quad
Considere el modelo de regresion simple con heteroscedasticidad:
\begin{equation}
  Y_i = \beta_1 + \beta_2 X_i + u_i,
  \qquad \E[u_i^2] = \sigma_i^2
  \label{eq:modelo}
\end{equation}
El estimador de Minimos Cuadrados Ponderados (MCP) con pesos $w_i = 1/\sigma_i^2$
es, segun la ecuacion (11.3.8) de Gujarati:
\begin{equation}
  \hat\beta_2^* = \frac{(\sw)(\swxy) - (\swx)(\swy)}{(\sw)(\swxx) - (\swx)^2}
  \label{eq:beta2_original}
\end{equation}
con varianza (ecuacion 11.3.9):
\begin{equation}
  \Var(\hat\beta_2^*) = \frac{\sw}{(\sw)(\swxx) - (\swx)^2}
  \label{eq:var_original}
\end{equation}

Se pide mostrar que estas expresiones pueden reescribirse como:
\begin{align}
  \hat\beta_2^* &= \frac{\sum w_i \ystar \xstar}{\sum w_i (\xstar)^2}
  \label{eq:beta2_nuevo}\\[6pt]
  \Var(\hat\beta_2^*) &= \frac{1}{\sum w_i (\xstar)^2}
  \label{eq:var_nuevo}
\end{align}
donde las desviaciones y medias ponderadas se definen como:
\begin{align}
  \Ybar &= \frac{\swy}{\sw}, \qquad \Xbar = \frac{\swx}{\sw}
  \label{eq:medias_pond}\\[6pt]
  \ystar &= Y_i - \Ybar, \qquad \xstar = X_i - \Xbar
  \label{eq:desviaciones}
\end{align}

\section{Lemas auxiliares}

Antes de la demostracion principal, establecemos tres lemas que simplifican el
algebra.

\begin{lema}
\label{lema1}
$\displaystyle\sum w_i \xstar = 0$
\end{lema}
\begin{proof}
\begin{equation}
  \sum w_i \xstar = \sum w_i(X_i - \Xbar)
               = \swx - \Xbar \sw
               = \swx - \frac{\swx}{\sw}\cdot\sw
               = \swx - \swx = 0 \qquad\square
\end{equation}
\end{proof}

\begin{lema}
\label{lema2}
$\displaystyle\sum w_i \ystar = 0$
\end{lema}
\begin{proof}
Analogamente al Lema~\ref{lema1}:
$\sum w_i \ystar = \swy - \Ybar\sw = \swy - \swy = 0$. \hfill$\square$
\end{proof}

\begin{lema}
\label{lema3}
$\displaystyle\sum w_i (\xstar)^2 = \swxx - \frac{(\swx)^2}{\sw}
 = \frac{(\sw)(\swxx) - (\swx)^2}{\sw}$
\end{lema}
\begin{proof}
\begin{align}
  \sum w_i(\xstar)^2
  &= \sum w_i(X_i - \Xbar)^2
   = \sum w_i\bigl(X_i^2 - 2X_i\Xbar + (\Xbar)^2\bigr)
  \notag\\
  &= \swxx - 2\Xbar\swx + (\Xbar)^2\sw
  \notag\\
  &= \swxx - 2\cdot\frac{\swx}{\sw}\cdot\swx + \frac{(\swx)^2}{(\sw)^2}\cdot\sw
  \notag\\
  &= \swxx - 2\frac{(\swx)^2}{\sw} + \frac{(\swx)^2}{\sw}
  \notag\\
  &= \swxx - \frac{(\swx)^2}{\sw}
   = \frac{(\sw)(\swxx) - (\swx)^2}{\sw}
  \label{eq:lema3_resultado}
  \qquad\square
\end{align}
\end{proof}

\section{Demostracion de \texorpdfstring{$\hat\beta_2^* = \sum w_i \ystar \xstar / \sum w_i (\xstar)^2$}{beta2*}}

\subsection{Desarrollo del numerador}

Expandimos $\sum w_i \ystar \xstar$ usando las definiciones~(\ref{eq:desviaciones}):

\begin{align}
  \sum w_i \ystar \xstar
  &= \sum w_i (Y_i - \Ybar)(X_i - \Xbar)
  \notag\\
  &= \sum w_i\bigl[Y_i X_i - Y_i\Xbar - X_i\Ybar + \Ybar\Xbar\bigr]
  \notag\\
  &= \swxy
   - \Xbar\swy
   - \Ybar\swx
   + \Ybar\Xbar\sw
  \label{eq:num_paso1}
\end{align}

Sustituyendo $\Xbar = \swx/\sw$ y $\Ybar = \swy/\sw$ en~(\ref{eq:num_paso1}):

\begin{align}
  &= \swxy
   - \frac{\swx}{\sw}\cdot\swy
   - \frac{\swy}{\sw}\cdot\swx
   + \frac{\swy}{\sw}\cdot\frac{\swx}{\sw}\cdot\sw
  \notag\\
  &= \swxy
   - \frac{(\swx)(\swy)}{\sw}
   - \frac{(\swy)(\swx)}{\sw}
   + \frac{(\swy)(\swx)}{\sw}
  \notag\\
  &= \swxy - \frac{(\swx)(\swy)}{\sw}
  \label{eq:num_resultado}
\end{align}

Multiplicando y dividiendo por $\sw$:

\begin{equation}
  \sum w_i \ystar \xstar
    = \frac{(\sw)(\swxy) - (\swx)(\swy)}{\sw}
  \label{eq:num_final}
\end{equation}

\subsection{Denominador}

Por el Lema~\ref{lema3}:
\begin{equation}
  \sum w_i(\xstar)^2
    = \frac{(\sw)(\swxx) - (\swx)^2}{\sw}
  \label{eq:den_final}
\end{equation}

\subsection{Cociente}

Dividiendo~(\ref{eq:num_final}) entre~(\ref{eq:den_final}):

\begin{equation}
  \frac{\sum w_i \ystar \xstar}{\sum w_i(\xstar)^2}
    = \frac{\displaystyle\frac{(\sw)(\swxy) - (\swx)(\swy)}{\sw}}
           {\displaystyle\frac{(\sw)(\swxx) - (\swx)^2}{\sw}}
    = \frac{(\sw)(\swxy) - (\swx)(\swy)}{(\sw)(\swxx) - (\swx)^2}
  \label{eq:cociente}
\end{equation}

El factor $1/\sw$ se cancela en numerador y denominador. La expresion resultante
es exactamente la ecuacion~(\ref{eq:beta2_original}). Por tanto:

\begin{equation}
  \boxed{
    \hat\beta_2^* = \frac{\sum w_i \ystar \xstar}{\sum w_i(\xstar)^2}
  }
  \label{eq:beta2_demostrado}
\end{equation}

\hfill$\square$

\section{Demostracion de \texorpdfstring{$\Var(\hat\beta_2^*) = 1/\sum w_i(\xstar)^2$}{var(beta2*)}}

\subsection{Punto de partida}

La varianza original~(\ref{eq:var_original}) es:
\begin{equation}
  \Var(\hat\beta_2^*) = \frac{\sw}{(\sw)(\swxx) - (\swx)^2}
  \label{eq:var_punto_partida}
\end{equation}

\subsection{Conexion con el denominador del Lema~\ref{lema3}}

Del Lema~\ref{lema3}:
\begin{equation}
  \sum w_i(\xstar)^2 = \frac{(\sw)(\swxx) - (\swx)^2}{\sw}
  \label{eq:lema3_rep}
\end{equation}

Despejando el denominador de~(\ref{eq:var_punto_partida}):
\begin{equation}
  (\sw)(\swxx) - (\swx)^2 = \sw \cdot \sum w_i(\xstar)^2
  \label{eq:despeje}
\end{equation}

\subsection{Sustitucion}

Sustituyendo~(\ref{eq:despeje}) en~(\ref{eq:var_punto_partida}):
\begin{equation}
  \Var(\hat\beta_2^*)
    = \frac{\sw}{\sw \cdot \sum w_i(\xstar)^2}
    = \frac{1}{\sum w_i(\xstar)^2}
  \label{eq:var_final}
\end{equation}

El factor $\sw$ se cancela en numerador y denominador:

\begin{equation}
  \boxed{
    \Var(\hat\beta_2^*) = \frac{1}{\sum w_i(\xstar)^2}
  }
  \label{eq:var_demostrada}
\end{equation}

\hfill$\square$

\section{Interpretacion de los resultados}

\subsection{Analogia con MCO ordinario}

Comparando los resultados obtenidos con los estimadores MCO ordinarios en
desviaciones:

\begin{center}
\small
\begin{tabular}{p{4.5cm} p{5cm} p{4.5cm}}
\toprule
\textbf{Cantidad}
  & \textbf{MCO ordinario}
  & \textbf{MCP (ponderado)} \\
\midrule
Pendiente
  & $\hat\beta_2 = \dfrac{\sum y_i x_i}{\sum x_i^2}$
  & $\hat\beta_2^* = \dfrac{\sum w_i \ystar \xstar}{\sum w_i (\xstar)^2}$ \\[10pt]
Varianza
  & $\Var(\hat\beta_2) = \dfrac{\sigma^2}{\sum x_i^2}$
  & $\Var(\hat\beta_2^*) = \dfrac{1}{\sum w_i (\xstar)^2}$ \\[10pt]
Medias usadas
  & $\bar Y = \sum Y_i/n$, $\bar X = \sum X_i/n$
  & $\Ybar = \sum w_i Y_i/\sum w_i$, $\Xbar = \sum w_i X_i/\sum w_i$ \\[6pt]
Desviaciones
  & $y_i = Y_i - \bar Y$, $x_i = X_i - \bar X$
  & $\ystar = Y_i - \Ybar$, $\xstar = X_i - \Xbar$ \\
\bottomrule
\end{tabular}
\end{center}

\vspace{6pt}

La estructura es completamente analoga: MCP es equivalente a MCO aplicado a
variables expresadas en desviaciones respecto a sus \emph{medias ponderadas},
donde las ponderaciones son $w_i = 1/\sigma_i^2$.

\subsection{Interpretacion de la varianza}

La varianza del estimador MCP es inversamente proporcional a la suma ponderada
de cuadrados de las desviaciones de $X$:
\begin{equation}
  \Var(\hat\beta_2^*) = \frac{1}{\sum w_i (\xstar)^2}
\end{equation}

Cuando las observaciones con mayor varianza del error (menor $w_i$) tienen
valores de $X$ alejados de la media, reciben menos peso, lo que es correcto
econometricamente: las observaciones ``ruidosas'' contribuyen menos a la
estimacion.

\subsection{Eficiencia del estimador MCP}

El estimador MCP es MELI (Mejor Estimador Lineal Insesgado) bajo heteroscedasticidad.
Esto se puede ver comparando su varianza con la del estimador MCO ordinario
aplicado al mismo modelo:

\begin{equation}
  \Var(\hat\beta_2^{\text{MCO}})
    = \frac{\sum w_i^2 (\xstar)^2}{\left[\sum w_i(\xstar)^2\right]^2}
    \cdot \frac{1}{w_i^2}
  \geq \frac{1}{\sum w_i(\xstar)^2}
    = \Var(\hat\beta_2^*)
  \label{eq:eficiencia}
\end{equation}

Por la desigualdad de Cauchy-Schwarz:
$\left[\sum w_i(\xstar)^2\right]^2 \leq \sum w_i(\xstar)^4 \cdot n$,
lo que confirma que MCP tiene menor o igual varianza que MCO bajo heteroscedasticidad.

\begin{clave}
  \textbf{Resultado fundamental.} El estimador MCP de la pendiente y su varianza:
  \[
    \hat\beta_2^* = \frac{\sum w_i \ystar \xstar}{\sum w_i (\xstar)^2},
    \qquad
    \Var(\hat\beta_2^*) = \frac{1}{\sum w_i (\xstar)^2}
  \]
  son algebraicamente identicos a las expresiones originales~(\ref{eq:beta2_original})
  y~(\ref{eq:var_original}). La representacion en desviaciones ponderadas hace
  explicita la analogia con MCO ordinario y muestra que MCP es equivalente a
  aplicar MCO con observaciones reescaladas por $\sqrt{w_i}$.
\end{clave}

\appendix

\section{Verificacion: propiedad de suma cero de las desviaciones ponderadas}

Una propiedad clave que subyace a todo el desarrollo es que las desviaciones
ponderadas suman cero:
\begin{align}
  \sum w_i \xstar &= 0 \quad \text{(Lema~\ref{lema1})} \label{eq:prop1}\\
  \sum w_i \ystar &= 0 \quad \text{(Lema~\ref{lema2})} \label{eq:prop2}
\end{align}

Estas propiedades son el analogo ponderado de la propiedad MCO ordinaria
$\sum x_i = \sum y_i = 0$. Garantizan que las condiciones normales del problema
MCP se reducen al sistema estudiado.

\section{Relacion entre MCP y MCO sobre datos transformados}

Si se define $Y_i^+ = \sqrt{w_i} Y_i$, $X_i^+ = \sqrt{w_i} X_i$ y
$\mathbf{1}_i^+ = \sqrt{w_i}$, entonces aplicar MCO ordinario al modelo
transformado:
\begin{equation}
  Y_i^+ = \beta_1 \mathbf{1}_i^+ + \beta_2 X_i^+ + u_i^+
\end{equation}
produce exactamente los estimadores MCP $\hat\beta_1^*$ y $\hat\beta_2^*$.
El error transformado $u_i^+ = u_i/\sqrt{\sigma_i^2}$ tiene varianza constante~1,
lo que restaura la homoscedasticidad y permite que MCO sea MELI sobre los
datos transformados.

\vspace{14pt}
\noindent\rule{\linewidth}{0.4pt}

\noindent{\small\textit{Nota.} Las ecuaciones~(\ref{eq:beta2_demostrado}) y~(\ref{eq:var_demostrada})
son las formas computacionalmente mas convenientes del estimador MCP. La
expresion en desviaciones ponderadas hace explicita la analogia estructural
con MCO y facilita la intuicion sobre como cada observacion contribuye
ponderadamente a la estimacion de la pendiente.}

\end{document}
